\documentclass[11pt]{beamer}
\usetheme{Warsaw}
\usepackage[utf8]{inputenc}
\usepackage[english]{babel}
\usepackage{amsmath}
\usepackage{amsfonts}
\usepackage{amssymb}
\usepackage{graphicx}
\usepackage{apacite}
\author{
	Kho, Lee
	\and
	Peeters, Jonas
}
\title{Sparse Principal Components Analysis}
% \setbeamercovered{transparent} 
% \setbeamertemplate{navigation symbols}{} 
% \logo{} 
\institute{\includegraphics{../Figures/nyu_long_black.png}} 
\date{\today} 
\subject{DS-GA 1013} 

\begin{document}

\begin{frame}
\titlepage
\end{frame}

\begin{frame}\frametitle{Introduction}
\begin{itemize}
    \item PCA is a popular dimension reduction tool with applications across many domains; however, one drawback of the technique is that PCs are usually a linear combination of \textit{all} variables and are thus difficult to interpret
    \begin{itemize}
        \item This is especially challenging in high-dimension, low sample size (HDLSS) data sets (i.e. "fat" matrices)
        \item PCA may also lead to inconsistent results in HDLSS data
    \end{itemize}
    \item \textbf{Sparse PCA} methods have been developed in an effort to compute more interpretable PCs with few nonzero loadings
    \item In our project, we implemented three of these sparse PCA methods and evaluated their results on breast tissue data set
\end{itemize}
\end{frame}

\begin{frame}\frametitle{State of the Art}
\begin{itemize}
    \item Simple thresholding, rotation and integer programming methods are sometimes used in practice to obtain sparse PCs, but these methods have been shown to produce non-optimal or even misleading results
    \item Other sparse PCA techniques such as \textbf{SCoT}, \textbf{SCoTLASS} and \textbf{SPCA} have been developed since the early 2000s that more effectively produce sparse PCs that maximize explained variance
    \begin{itemize}
    \item Further adaptations and improvements have been made on these sparse PCA techniques, but they continue to be frequently cited in related literature as core methods
    \end{itemize}
\end{itemize}
\end{frame}

\begin{frame}\frametitle{Methodology: Selected Sparse PCA Methods}
We implemented the \textbf{SCoT}, \textbf{SCoTLASS} and \textbf{SPCA} methods and on a breast tissue measurements data set with n = 106 and p = 9\\
\begin{itemize}
    \item \textbf{SCoT}: This method maximizes both the regular PCA objective and the \textit{varimax} criterion, which is a sparsity criterion that maximizes the variance of the squared loadings
    \item \textbf{SCoTLASS}: SCoTLASS induces sparse loadings by maximizing the regular PCA subject to an \textit{L1-norm} constraint
    \item \textbf{SPCA}: SPCA re-configures the sparse PCA problem into a \textit{regression problem} and uses an \textit{elastic net} constraint to efficiently produce sparse loadings 
\end{itemize}
To optimize the SCoT and SCoTLASS objective functions, we used projected gradient ascent
\begin{itemize}
    \item Due to the non-differentiable nature of the L1-norm, we approximated the norm with a smooth function (this procedure was outlined by Trendafilov (2004))
\end{itemize}
\end{frame}

\begin{frame}\frametitle{Methodology: Evaluation}
We evaluated these three sparse PCA methods based on the following three criteria:
\begin{itemize}
    \item \textit{Sparsity}: We looked at sparsity in terms of the \textit{gini coefficient}, which is a metric for weight distribution across loadings
        \begin{itemize}
            \item The Gini coefficient can take on a value between 0 and 1 inclusive, with 0 indicating minimal sparsity (i.e. all loadings have equal weight) and 1 indicating maximum sparsity (i.e. all weight on one loading)
        \end{itemize}
    \item \textit{Explained variance}: We calculated \textit{percentage of explained variance (PEV)} of the sparse PCs that resulted from each method
    \item \textit{Efficiency and scalability}: For each method, we measured run time on the breast tissue data set assessed the method's suitability for higher dimensional data
\end{itemize}
\end{frame}

\begin{frame}\frametitle{Results}
\begin{itemize}
\item SPCA proved to be the most efficient and scalable method that we explored
\begin{itemize}
    \item Both SCoT and SCoTLASS require hyperparameter selection, which turned out to be fairly tricky and time consuming; in addition, the non-convex nature of the SCoTLASS objective can result in many local maxima, and thus several random initialization points may be necessary to find the global maximum
\end{itemize}
\end{itemize}
\end{frame}

\begin{frame}\frametitle{Discussion}
\begin{itemize}
\item 
\end{itemize}
\end{frame}

\end{document}